% !TEX encoding = UTF-8 Unicode
\documentclass[11pt,a4paper]{article}

\usepackage[T1]{fontenc}
\usepackage[utf8]{inputenc}
\usepackage[french]{babel}
\usepackage{microtype}
\usepackage{geometry}
\geometry{margin=2.5cm}
\usepackage{booktabs}
\usepackage{array}
\usepackage{hyperref}
\usepackage{xcolor}
\usepackage{enumitem}
\usepackage{titlesec}

\hypersetup{
  colorlinks=true,
  linkcolor=blue!60!black,
  urlcolor=blue!60!black,
  citecolor=blue!60!black
}

\title{Rapport d'assurance qualité logicielle --- Module frontend\\[0.5em]
\large Analyse statique et suivi des indicateurs (SonarQube)}
\author{Équipe de 3 étudiants --- \textit{[Noms à compléter]}\\
\textit{[Module d'assurance qualité logicielle --- à compléter]}\\
\textit{[Enseignant référent --- à compléter]}}
\date{27 mars 2026}

\begin{document}

\maketitle

\begin{abstract}
\noindent
Ce rapport synthétise les résultats d'analyse statique du frontend \texttt{ecommerce-frontend} à partir d'exports CSV SonarQube (métriques instantanées, historique, issues). Il interprète les indicateurs de qualité, identifie les risques et propose un plan d'actions pour une équipe de trois étudiants.
\end{abstract}

\vspace{0.8em}
\noindent\textbf{Contexte :} travail réalisé dans le cadre d'un module d'\textbf{assurance de la qualité logicielle} à l'université.

\noindent\textbf{Périmètre :} application \textbf{e-commerce} --- partie frontend (\texttt{ecommerce-frontend}).

\noindent\textbf{Sources de données :} exports CSV générés à partir du serveur SonarQube :
\begin{itemize}[noitemsep]
  \item \texttt{rapport-sonar-metrics.csv} --- instantané des métriques projet ;
  \item \texttt{rapport-sonar-metrics-history.csv} --- historique des métriques suivies dans le temps ;
  \item \texttt{rapport-sonar-issues.csv} --- détail des anomalies (issues) remontées par l'analyseur.
\end{itemize}

\noindent\textbf{Date de référence des métriques :} 27 mars 2026 (dernier état reflété dans les CSV).

\section{Introduction}

\subsection{Objectif du rapport}

Ce document présente une \textbf{synthèse académique} des résultats d'analyse de code du frontend : taille, complexité, dette technique, couverture de tests, sécurité, fiabilité, et liste des problèmes détectés. L'objectif est de \textbf{documenter l'état de la qualité} du logiciel, d'\textbf{interpréter} les indicateurs et de \textbf{proposer des pistes d'amélioration} alignées avec les bonnes pratiques d'ingénierie logicielle.

\subsection{Méthodologie}

L'outil \textbf{SonarQube} effectue une \textbf{analyse statique} du code source (TypeScript/JavaScript, HTML/CSS selon configuration). Les métriques et issues sont agrégées au niveau du composant projet. Les fichiers CSV utilisés ici constituent une \textbf{trace reproductible} des résultats au moment de l'analyse.

\textit{Note :} le fichier \texttt{rapport-sonar-issues.csv} contient également un \textbf{historique d'issues} (dont beaucoup \textbf{fermées} après correction ou reconfiguration), ce qui explique un volume de lignes bien supérieur au nombre d'issues \textbf{ouvertes} actuellement pris en compte dans le tableau de bord (voir section~\ref{sec:issues}).

\section{Synthèse exécutive}

\begin{table}[h]
  \centering
  \begin{tabular}{@{}p{3.8cm}p{10cm}@{}}
    \toprule
    \textbf{Domaine} & \textbf{Synthèse} \\
    \midrule
    Taille & 3\,205 lignes de code (NCLOC), complexité cyclomatique totale \textbf{465}. \\
    Fiabilité & \textbf{1 bug} ouvert ; note de fiabilité \textbf{B} (rating 2.0). \\
    Sécurité & \textbf{1 vulnérabilité} ouverte, \textbf{1 point sensible} (security hotspot) ; \textbf{note de sécurité E} (rating 5.0) --- priorité à traiter. \\
    Maintenabilité & \textbf{30 code smells} ouverts ; dette technique estimée \textbf{142 min}, ratio \textbf{0,1\,\%}. \\
    Tests & Couverture \textbf{53,8\,\%} ; \textbf{890} lignes à couvrir, \textbf{445} non couvertes. \\
    Duplication & \textbf{0\,\%} de lignes dupliquées (densité), \textbf{0} bloc dupliqué. \\
    Documentation inline & Densité de commentaires \textbf{0,8\,\%} (faible). \\
    \bottomrule
  \end{tabular}
\end{table}

\textbf{Conclusion provisoire :} le projet présente une \textbf{base de code de taille raisonnable}, une \textbf{duplication maîtrisée}, et une \textbf{couverture de tests en progression} mais encore \textbf{sous les seuils souvent visés en industrie ($\geq 80\,\%$)}. En revanche, la \textbf{sécurité} et un \textbf{bug} restent des \textbf{points d'attention immédiats}.

\section{Analyse détaillée des métriques instantanées}

Les valeurs ci-dessous proviennent de \texttt{rapport-sonar-metrics.csv}.

\subsection{Taille et structure}

\begin{table}[h]
  \centering
  \begin{tabular}{@{}llp{7cm}@{}}
    \toprule
    \textbf{Métrique} & \textbf{Valeur} & \textbf{Interprétation} \\
    \midrule
    \texttt{ncloc} & 3\,205 & Volume fonctionnel du code (hors commentaires/blancs significatifs). \\
    \texttt{complexity} & 465 & Complexité cyclomatique cumulée ; à surveiller sur les modules les plus denses. \\
    \texttt{comment\_lines\_density} & 0,8\,\% & Peu de commentaires explicatifs ; la lisibilité repose surtout sur le nommage et la structure. \\
    \bottomrule
  \end{tabular}
\end{table}

\subsection{Qualité intrinsèque (bugs, smells, dette)}

\begin{table}[h]
  \centering
  \begin{tabular}{@{}llp{6.5cm}@{}}
    \toprule
    \textbf{Métrique} & \textbf{Valeur} & \textbf{Interprétation} \\
    \midrule
    \texttt{bugs} & 1 & Au moins un défaut potentiel de comportement à corriger. \\
    \texttt{code\_smells} & 30 & Améliorations de style, conception ou bonnes pratiques. \\
    \texttt{vulnerabilities} & 1 & Risque de sécurité identifié par les règles SonarQube. \\
    \texttt{violations} & 32 & Total cohérent avec la somme des problèmes ouverts suivis (voir issues). \\
    \texttt{sqale\_index} & 142 min & Effort estimé pour « rembourser » la dette technique. \\
    \texttt{sqale\_debt\_ratio} & 0,1\,\% & Dette faible par rapport à la taille du code. \\
    \texttt{reliability\_rating} & 2,0 & \textbf{B} --- bon niveau global de fiabilité, sous réserve de corriger le bug ouvert. \\
    \texttt{security\_rating} & 5,0 & \textbf{E} --- niveau le plus critique sur l'échelle Sonar ; \textbf{action prioritaire} sur la sécurité. \\
    \texttt{security\_hotspots} & 1 & Zone à revue manuelle pour confirmer ou infirmer un risque. \\
    \bottomrule
  \end{tabular}
\end{table}

\subsection{Tests et couverture}

\begin{table}[h]
  \centering
  \begin{tabular}{@{}llp{7cm}@{}}
    \toprule
    \textbf{Métrique} & \textbf{Valeur} & \textbf{Interprétation} \\
    \midrule
    \texttt{coverage} & 53,8\,\% & Environ la moitié du code « couvert » par les tests automatisés. \\
    \texttt{lines\_to\_cover} & 890 & Volume de code concerné par la mesure de couverture. \\
    \texttt{uncovered\_lines} & 445 & Cible prioritaire pour étendre les tests. \\
    \bottomrule
  \end{tabular}
\end{table}

\subsection{Duplication}

\begin{table}[h]
  \centering
  \begin{tabular}{@{}ll@{}}
    \toprule
    \textbf{Métrique} & \textbf{Valeur} \\
    \midrule
    \texttt{duplicated\_lines\_density} & 0,0\,\% \\
    \texttt{duplicated\_blocks} & 0 \\
    \bottomrule
  \end{tabular}
\end{table}

\textbf{Interprétation :} absence de duplication détectée aux seuils de l'outil --- point \textbf{positif} pour la maintenabilité.

\section{Analyse des issues (\texttt{rapport-sonar-issues.csv})}
\label{sec:issues}

\subsection{Volume global}

\begin{itemize}
  \item \textbf{Total d'enregistrements d'issues} dans le CSV : \textbf{3\,933} (hors ligne d'en-tête), incluant l'historique (issues \textbf{CLOSED} et \textbf{OPEN}).
  \item \textbf{Issues actuellement ouvertes (OPEN)} : \textbf{32}, en cohérence avec la métrique \texttt{violations} $= 32$.
\end{itemize}

\subsection{Répartition des 32 issues ouvertes par sévérité}

\begin{table}[h]
  \centering
  \begin{tabular}{@{}lr@{}}
    \toprule
    \textbf{Sévérité} & \textbf{Nombre} \\
    \midrule
    \textbf{BLOCKER} & 1 \\
    \textbf{MAJOR} & 10 \\
    \textbf{MINOR} & 21 \\
    \bottomrule
  \end{tabular}
\end{table}

\subsection{Répartition par type SonarQube}

\begin{table}[h]
  \centering
  \begin{tabular}{@{}lr@{}}
    \toprule
    \textbf{Type} & \textbf{Nombre} \\
    \midrule
    CODE\_SMELL & 30 \\
    VULNERABILITY & 1 \\
    BUG & 1 \\
    \bottomrule
  \end{tabular}
\end{table}

\subsection{Règles les plus fréquentes (issues ouvertes)}

Les règles suivantes apparaissent le plus souvent parmi les issues \textbf{OPEN} :

\begin{table}[h]
  \centering
  \small
  \begin{tabular}{@{}p{3.2cm}cp{7.5cm}@{}}
    \toprule
    \textbf{Règle} & \textbf{Occ.} & \textbf{Thème typique} \\
    \midrule
    \texttt{typescript:S6759} & 11 & Props de composants React à marquer en lecture seule (\texttt{readonly}). \\
    \texttt{typescript:S7763} & 7 & Préférer \texttt{export\ldots from} pour les réexports de modules. \\
    \texttt{typescript:S3358} & 3 & Simplifier des expressions conditionnelles imbriquées. \\
    \texttt{typescript:S6557} & 3 & Ajustements TypeScript sur les promesses / async. \\
    \texttt{typescript:S1128} & 2 & Imports inutilisés à supprimer. \\
    \texttt{secrets:S6702} & 1 & \textbf{Revues liées aux secrets} (sensibilité sécurité --- à traiter en priorité avec la vulnérabilité). \\
    \bottomrule
  \end{tabular}
\end{table}

\textit{Les libellés exacts des règles sont disponibles dans l'interface SonarQube ; le CSV liste la clé de règle et le message par issue.}

\subsection{Issues historiques (fermées)}

La majorité des lignes du CSV correspondent à des issues \textbf{CLOSED} (anomalies corrigées ou analyses passées incluant notamment des artefacts sous \texttt{coverage/}). Cela illustre l'\textbf{évolution} du projet et la \textbf{capacité à résorber} de la dette, mais ne reflète pas l'état actuel du tableau de bord pour la \textbf{qualité du code applicatif} seul.

\section{Évolution temporelle (\texttt{rapport-sonar-metrics-history.csv})}

L'historique disponible couvre une fenêtre récente (26--27 mars 2026) pour les métriques : \texttt{bugs}, \texttt{code\_smells}, \texttt{coverage}, \texttt{vulnerabilities}.

\subsection{Faits marquants}

\begin{itemize}
  \item \textbf{Bugs :} passage transitoire à \textbf{25} sur un instantané du 26/03, puis \textbf{retour à 1} --- indique une \textbf{instabilité ponctuelle} (branche, périmètre d'analyse ou fusion) puis stabilisation.
  \item \textbf{Code smells :} pic très élevé (\textbf{3896}) sur un point du 26/03, puis valeurs basses (\textbf{26--30}) ensuite --- \textbf{cohérent avec une analyse ayant inclus des fichiers non représentatifs} (ex.\ générés) avant filtrage ; les valeurs récentes (\textbf{30}) sont alignées avec le tableau de bord actuel.
  \item \textbf{Couverture :} de \textbf{0\,\%} à \textbf{53,8\,\%} sur la période --- \textbf{progression nette} une fois les rapports de couverture pris en compte par SonarQube.
  \item \textbf{Vulnérabilités :} stable à \textbf{1} sur la période --- \textbf{non résolue} à ce stade.
\end{itemize}

\subsection{Lecture pédagogique}

Ce type de courbe montre l'importance de \textbf{conditions d'analyse stables} (même branche, mêmes exclusions, rapport LCOV à jour) pour comparer les versions dans le temps.

\section{Risques et priorités}

\begin{enumerate}[label=\arabic*.]
  \item \textbf{Sécurité (note E, vulnérabilité + hotspot + règle \texttt{secrets:S6702})} --- analyser en équipe, corriger ou justifier, puis \textbf{revue de code} ciblée.
  \item \textbf{Bug ouvert} --- reproduction, test de non-régression, correction.
  \item \textbf{Issue BLOCKER} --- traitement \textbf{immédiat} selon la politique qualité du module.
  \item \textbf{Couverture $\sim$54\,\%} --- planifier des tests sur les chemins critiques (authentification, panier, commande, erreurs API).
  \item \textbf{Dette « smells » TypeScript/React} --- corrections \textbf{par lots} (réexports, \texttt{readonly}, nettoyage d'imports) pour réduire le bruit dans SonarQube.
\end{enumerate}

\section{Plan d'actions recommandé (équipe de 3 étudiants)}

\begin{table}[h]
  \centering
  \begin{tabular}{@{}llp{6.5cm}@{}}
    \toprule
    \textbf{Priorité} & \textbf{Action} & \textbf{Livrable} \\
    \midrule
    P0 & Atelier sécurité : vulnérabilité, hotspot, règle secrets & Correctif ou rapport de justification + nouvelle analyse Sonar \\
    P0 & Corriger le bug et l'issue BLOCKER & PR + tests \\
    P1 & Uniformiser les réexports (\texttt{S7763}) et props (\texttt{S6759}) & PRs thématiques par dossier (\texttt{api}, \texttt{components}, \ldots) \\
    P2 & Augmenter la couverture (cible intermédiaire \textbf{65\,\%} puis \textbf{80\,\%}) & Rapport de couverture + nouveaux tests \\
    P3 & Documenter les modules complexes (commentaires ciblés) & Révision \texttt{comment\_lines\_density} \\
    \bottomrule
  \end{tabular}
\end{table}

Répartition possible : \textbf{étudiant A} --- sécurité et secrets ; \textbf{étudiant B} --- bugs et tests ; \textbf{étudiant C} --- smells TypeScript/React et duplication de configuration d'analyse.

\section{Conclusion}

Les CSV SonarQube décrivent un \textbf{frontend structuré} (faible duplication, dette relativement faible) avec une \textbf{couverture de tests en amélioration continue} mais encore \textbf{insuffisante} pour les standards industriels exigeants. Les \textbf{points critiques} sont \textbf{sécurité} (note E, vulnérabilité persistante) et la \textbf{résolution du bug} et du \textbf{BLOCKER}. La poursuite des \textbf{bonnes pratiques TypeScript/React} (issues \texttt{S6759}, \texttt{S7763}, etc.) permettra de stabiliser la qualité \textbf{maintenabilité} sans alourdir excessivement la charge.

Ce rapport peut être \textbf{joint} à une présentation orale en citant explicitement les \textbf{fichiers sources CSV} et la \textbf{date d'analyse}, et en rappelant que la qualité est un \textbf{processus itératif} mesuré par des outils mais \textbf{validé} par des revues humaines et des tests.

\section*{Annexes}
\addcontentsline{toc}{section}{Annexes}

\subsection*{A. Correspondance indicative des ratings SonarQube (rappel)}

\begin{itemize}[noitemsep]
  \item \textbf{Fiabilité / sécurité / maintenabilité :} échelle \textbf{1 (A)} à \textbf{5 (E)} ; \textbf{1} = meilleur niveau selon les règles du produit.
\end{itemize}

\subsection*{B. Fichiers de données}

\begin{itemize}[noitemsep]
  \item \texttt{rapport-sonar-metrics.csv}
  \item \texttt{rapport-sonar-metrics-history.csv}
  \item \texttt{rapport-sonar-issues.csv}
\end{itemize}

\subsection*{C. Références bibliographiques indicatives}

\begin{itemize}[noitemsep]
  \item ISO/IEC 25010 --- \textit{Systems and software Quality Requirements and Evaluation (SQuaRE)} --- modèle de qualité du produit logiciel.
  \item Documentation SonarQube --- \textit{Metrics definitions} et \textit{Security} (selon version utilisée en cours).
  \item Cours d'assurance qualité logicielle --- méthodes de revue statique et dynamique.
\end{itemize}

\vspace{1.5em}
\noindent\textit{Document rédigé pour un usage pédagogique --- à compléter par les noms des étudiants, l'intitulé exact du cours, l'enseignant référent et la date de remise.}

\end{document}
